\documentclass[11pt]{article}
\usepackage[T1]{fontenc}
\usepackage{lmodern,amsmath,amssymb,amsthm,mathtools,longtable,array}
\usepackage[margin=1.05in]{geometry}
\usepackage[hidelinks]{hyperref}
\newtheorem{theorem}{Theorem}[section]
\newtheorem{lemma}[theorem]{Lemma}
\newtheorem{proposition}[theorem]{Proposition}
\newtheorem{corollary}[theorem]{Corollary}
\theoremstyle{remark}\newtheorem{remark}[theorem]{Remark}
\newcommand{\R}{\mathbb R}\newcommand{\C}{\mathbb C}
\newcommand{\Z}{\mathbb Z}\newcommand{\T}{\mathbb T}
\newcommand{\norm}[1]{\left\lVert#1\right\rVert}
\newcommand{\code}[1]{{\small\nolinkurl{#1}}}
\DeclareMathOperator{\Dom}{Dom}
\title{Fourier covariance and positivity on the three-torus}
\author{Project Stardust}
\date{Reviewed editorial draft, October 29, 2026}
\begin{document}
\maketitle
\begin{abstract}
For normalized even Fourier symbols on the three-torus, we give
complete proofs connecting anchored covariance, conditional
negativity, pointwise iterated Gamma forms, finite heat-multiplier
matrices and positivity on the full nonnegative Haar-$L^2$ cone.
The passage from finite Fourier squares to arbitrary nonnegative
$L^2$ functions is proved explicitly. Coordinatewise power symbols
have the sharp positivity condition that every positive-weight
exponent is at most two; inactive coordinates impose no restriction.
We also construct the unique actual torus probability laws, their convolution kernels and countable-step chains. A genuine Poisson mixture has the exact transformed symbol, a persistent no-jump atom and a short-parameter jump-covariance identity. We identify the exact closed generator domain and explain why
semigroup existence alone supplies no positivity. A bounded rank-two
multiplier has an explicit nonnegative test sent outside the Haar cone
at every positive time. Exact integrated Gamma dissipation and sharp
all-level decay are also proved, including non-Markov power orders. A finite-time
increment kernel remains positive at every time, yet an exact
two-frequency quadratic-heat example shows that it need not increase
in Loewner order. The main equivalence is classical Schoenberg
theory, organized here around a frozen Lean development. No new
priority claim, completed pointwise Gamma domain or continuous-time
path construction is asserted.
\end{abstract}

\section{The question and the domains}
Let $\T^3=(\R/\Z)^3$, with normalized Haar measure $m$, and put
$\Lambda=\Z^3$. For $k\in\Lambda$ write
\[
 e_k(x)=\exp(2\pi i k\cdot x).
\]
The finite Fourier algebra $\mathcal P$ consists of the finite sums
$f=\sum_k\widehat f(k)e_k$. Its product is pointwise multiplication,
$e_ke_l=e_{k+l}$, and complex conjugation sends $e_k$ to $e_{-k}$.
The exponentials form an orthonormal basis of $L^2(\T^3,m;\C)$.
We use the unit-period normalization throughout. In particular the
rates below are specified on integer frequencies; a factor $2\pi$
is not silently inserted into their definition.

For a real symbol $\psi:\Lambda\to\R$, the algebraic operator
$L_\psi$ is defined on $\mathcal P$ by
\begin{equation}\label{eq:L}
 L_\psi e_k=-\psi(k)e_k.
\end{equation}
Until a closed domain is specified, this is only a finite Fourier
operator. The sesquilinear iterates used in the formal library are
\begin{align}\label{eq:gamma}
 \Gamma_0(f,g)&=\overline f g,\notag\\
 \Gamma_{n+1}(f,g)&=\tfrac12\bigl(
 L_\psi\Gamma_n(f,g)-\Gamma_n(L_\psi f,g)
                         -\Gamma_n(f,L_\psi g)\bigr).
\end{align}
They are conjugate linear in $f$ and linear in $g$. For an even
real symbol, $L_\psi\overline f=\overline{L_\psi f}$, and
$\Gamma_1$ is the usual complex carré du champ. We keep the
recursion itself as the definition when discussing a noneven
symbol in Section~\ref{sec:normalization}.

Define the anchored kernel
\begin{equation}\label{eq:K}
 K_\psi(k,l)=\frac{\psi(k)+\psi(l)-\psi(l-k)}2.
\end{equation}
Saying that $K_\psi$ is positive semidefinite means that every
finite matrix $(K_\psi(k_i,k_j))_{i,j=1}^n$ is real symmetric and
has nonnegative quadratic form. Repetitions of the sampled
frequencies are allowed. Conditional negativity means that
\begin{equation}\label{eq:CND}
 \sum_{i,j=1}^n c_i c_j\psi(k_j-k_i)\leq0
 \quad\hbox{whenever }c_i\in\R,\quad\sum_i c_i=0.
\end{equation}
Symmetry is stated separately; condition~\eqref{eq:CND} by itself
does not test the odd part of a real symbol.

The physical nonnegative cone in complex $L^2$ is
\[
 \mathcal C=\{u\in L^2(\T^3,m;\C):
          \operatorname{Im}u=0,\ \operatorname{Re}u\geq0
          \text{ almost everywhere}\}.
\]
Thus an inequality about Fourier coefficients is not the
definition of positivity used here.

\begin{theorem}\label{thm:main}
Suppose $\psi:\Lambda\to[0,\infty)$ is even and $\psi(0)=0$.
Let $S_t$ on $L^2$ be the Fourier multiplier
$\widehat{S_tu}(k)=e^{-t\psi(k)}\widehat u(k)$, for $t\geq0$.
The following conditions are equivalent:
\begin{enumerate}
\item $K_\psi$ is positive semidefinite on every finite sample;
\item $\psi$ satisfies~\eqref{eq:CND};
\item $\Gamma_1(f,f)(0)\geq0$ for every $f\in\mathcal P$;
\item $(e^{-t\psi(k_j-k_i)})_{i,j}$ is positive semidefinite
      for every finite sample and every $t\geq0$;
\item $S_t\mathcal C\subseteq\mathcal C$ for every $t\geq0$.
\end{enumerate}
Under these conditions, $\Gamma_n(f,f)(x)$ is real and
nonnegative for every integer $n\geq0$, $f\in\mathcal P$ and
$x\in\T^3$. The semigroup is strongly continuous, self-adjoint,
contractive and constant preserving. Its exact generator domain is
\begin{equation}\label{eq:domain}
 \Dom(A)=\left\{u\in L^2:
                  \sum_k\psi(k)^2|\widehat u(k)|^2<\infty\right\},
 \qquad\widehat{Au}(k)=-\psi(k)\widehat u(k).
\end{equation}
Finite Fourier polynomials form a core for this closed generator.
\end{theorem}

The finite equivalence in this statement is Schoenberg's
classical mechanism~\cite{schoenberg}. Its value here is the
complete connection to the specified physical $L^2$ realization,
including the square-cone closure and the exact operator domain.
All proofs are given below. Positivity of each finite matrix is
not used as a substitute for proving positivity on $\mathcal C$.

\section{Anchoring a conditionally negative kernel}
The zero frequency is needed as an extra sample even when it
does not occur in the polynomial of interest.

\begin{lemma}\label{lem:anchor}
Let $\psi$ be real, even and normalized by $\psi(0)=0$.
Then $K_\psi$ is positive semidefinite on finite samples if and
only if~\eqref{eq:CND} holds.
\end{lemma}
\begin{proof}
Evenness makes $K_\psi$ symmetric. For any coefficients $b_i$,
write $b=\sum_i b_i$. Expanding~\eqref{eq:K} gives
\begin{equation}\label{eq:anchoridentity}
 \sum_{i,j}b_i b_j K_\psi(k_i,k_j)
 =b\sum_i b_i\psi(k_i)
       -\frac12\sum_{i,j}b_i b_j\psi(k_j-k_i).
\end{equation}
For $b=0$, positivity of the left side yields conditional
negativity.

Conversely, augment the sample by $k_0=0$ and give it coefficient
$b_0=-\sum_{i=1}^n b_i$. The augmented coefficients sum to zero.
The diagonal term at the anchor vanishes, and its two cross terms
are $-2b\sum_i b_i\psi(k_i)$, by evenness. Therefore the
conditionally negative quadratic form of the augmented sample
is exactly $-2$ times the left side of~\eqref{eq:anchoridentity}.
Its nonpositivity proves the desired matrix positivity for
arbitrary $b_i$.
\end{proof}

A real positive semidefinite matrix $M$ also has nonnegative
complex energy: if $z=a+ib$, symmetry gives
\[
 \sum_{i,j}\overline z_i z_j M_{ij}
       =a^{\mathsf T}Ma+b^{\mathsf T}Mb\geq0.
\]
This observation permits complex Fourier polynomials in all
subsequent pointwise tests without changing the real definition
of conditional negativity.

\section{The entire hierarchy on the finite algebra}
For a single pair of modes, recursion~\eqref{eq:gamma} gives
\begin{equation}\label{eq:modegamma}
 \Gamma_n(e_k,e_l)=K_\psi(k,l)^n e_{l-k}.
\end{equation}
Indeed the formula is true for $n=0$; the next recursion multiplies
its coefficient by
$(-\psi(l-k)+\psi(k)+\psi(l))/2$. The symbol is real, so
conjugating the coefficient $-\psi(k)$ in the first argument
does not change it. Finite sesquilinearity now yields
\begin{equation}\label{eq:gammaenergy}
 \Gamma_n(f,f)(x)=\sum_{i,j}
  \overline{c_i e_{k_i}(x)}c_j e_{k_j}(x)
                      K_\psi(k_i,k_j)^n,
 \quad f=\sum_i c_i e_{k_i}.
\end{equation}
Here $M^{\circ n}$ denotes the entrywise power, not the matrix
power $M^n$.

\begin{lemma}\label{lem:schur}
If $M$ and $N$ are real positive semidefinite finite matrices,
then their entrywise product $M\circ N$ is positive semidefinite.
Consequently $M^{\circ n}$ is positive semidefinite for every
integer $n\geq0$.
\end{lemma}
\begin{proof}
Choose finite-dimensional Gram representations
$M_{ij}=\langle u_i,u_j\rangle$ and
$N_{ij}=\langle v_i,v_j\rangle$, for example by orthogonal
diagonalization and square roots of their nonnegative eigenvalues.
Then $(M\circ N)_{ij}$ is the Gram matrix of $u_i\otimes v_i$
in the finite tensor-product Euclidean space. Iteration proves
the power assertion, starting from the all-ones Gram matrix
$M^{\circ0}$.
\end{proof}

\begin{proposition}\label{prop:gammapositive}
For a real even normalized symbol, the following are equivalent:
$K_\psi$ is positive semidefinite; $\Gamma_1(f,f)(0)\geq0$
for every finite Fourier polynomial; and all the values
$\Gamma_n(f,f)(x)$ are real and nonnegative.
\end{proposition}
\begin{proof}
If the kernel is positive semidefinite, apply
Lemma~\ref{lem:schur} to~\eqref{eq:gammaenergy}, using the complex
energy observation after Lemma~\ref{lem:anchor}. This proves the
full hierarchy. The hierarchy implies its first origin test.
For the reverse implication, test
$f=\sum_i c_i e_{k_i}$ at $x=0$ with arbitrary real $c_i$.
Since $e_{k_i}(0)=1$, equation~\eqref{eq:gammaenergy} at $n=1$
is exactly the quadratic form of $K_\psi$ on that sample.
Repeated modes cause no difficulty: the formula is valid for
every finite presentation of $f$.
\end{proof}

This result is positivity on $\mathcal P$, not an assertion that
the products in~\eqref{eq:gamma} extend to all of $\Dom(A)$.
The exact generator domain alone does not control the products
or pointwise representatives required there.

\section{Normalization and real-symbol pitfalls}\label{sec:normalization}
The evenness assumption is structural. It must not be dropped
merely because a scalar quadratic form is called negative.

\begin{proposition}\label{prop:forces-even}
Let $\psi:\Lambda\to\R$ satisfy $\psi(0)=0$, without initially
assuming evenness. If
$\operatorname{Re}\Gamma_1(f,f)(0)\geq0$ for every
$f\in\mathcal P$, then $\psi$ is even. It follows that the
normalized first-origin-test criterion is equivalent to
conditional negativity together with evenness.
\end{proposition}
\begin{proof}
Fix $k$ and test $f=s+e_k$ with $s\in\R$. Formula
\eqref{eq:gammaenergy}, which follows from the recursion even
without evenness, gives
\[
 \operatorname{Re}\Gamma_1(s+e_k,s+e_k)(0)
 =\psi(k)+\frac{s}{2}\bigl(\psi(k)-\psi(-k)\bigr).
\]
An affine real function nonnegative for every $s\in\R$ has zero
slope. Hence $\psi(k)=\psi(-k)$. Apply
Proposition~\ref{prop:gammapositive} and Lemma~\ref{lem:anchor}
to finish both directions of the criterion.
\end{proof}

For an explicit failure, take $\psi(k)=k_1$. For every real
finite vector $c$, the matrix $\psi(k_j-k_i)$ is antisymmetric,
so its quadratic form is zero. In particular it satisfies the
bare inequality~\eqref{eq:CND}. Yet for $k=(1,0,0)$ and $s=-2$
the displayed first-Gamma value equals $1-2=-1$. This symbol
is not nonnegative and does not fall under
Theorem~\ref{thm:main}; it isolates the missing symmetry
assumption rather than contradicting that theorem.

Without normalization, positivity of only the real part can
also miss an imaginary antisymmetric contribution to
\eqref{eq:gammaenergy}. Testing $f=e_k+ie_l$ gives
\[
 \operatorname{Im}\Gamma_1(f,f)(0)
                   =K_\psi(k,l)-K_\psi(l,k).
\]
Thus the library's unnormalized complex criterion keeps reality
as well as nonnegativity. The main theorem avoids this ambiguity
by stating normalization and evenness at the outset.

\section{Finite Schoenberg positivity}
We next pass from the infinitesimal symbol to finite-time
multiplier matrices. This step is still a finite matrix argument.

\begin{proposition}\label{prop:schoenberg}
For real even $\psi$ with $\psi(0)=0$, conditional negativity
is equivalent to positivity of all finite matrices
\begin{equation}\label{eq:heatmatrix}
 H_t(i,j)=\exp\bigl(-t\psi(k_j-k_i)\bigr),\qquad t\geq0.
\end{equation}
\end{proposition}
\begin{proof}
By Lemma~\ref{lem:anchor}, assume the anchored matrix
$K=(K_\psi(k_i,k_j))$ is positive semidefinite. Its entrywise
exponential
\[
 \exp^{\circ}(2tK)=\sum_{n=0}^\infty
                         \frac{(2t)^n}{n!}K^{\circ n}
\]
is positive semidefinite. Each partial sum has that property by
Lemma~\ref{lem:schur}, and the entrywise convergent finite-matrix
limit preserves every quadratic-form inequality and symmetry.
Set $d_i=e^{-t\psi(k_i)}$. The identity
\[
 H_t(i,j)=d_i d_j\exp(2tK_{ij})
\]
shows that $H_t$ is the congruence of this positive matrix by the
real diagonal matrix with diagonal $d_i$. Hence it is positive
semidefinite.

Conversely, take arbitrary real $c_i$ with $\sum_i c_i=0$.
The finite sum
\[
 E(t)=\sum_{i,j}c_i c_j e^{-t\psi(k_j-k_i)}
\]
is nonnegative for $t\geq0$ and satisfies $E(0)=0$. Its right
derivative at zero is nonnegative. Differentiating its finite
number of terms gives
\[
 0\leq E'(0)=-\sum_{i,j}c_i c_j\psi(k_j-k_i),
\]
which is conditional negativity.
\end{proof}

No growth condition on $\psi$ is required for this finite
argument. For the physical $L^2$ evolution, nonnegativity of the
rates ensures bounded multipliers. Under the main normalized
even conditions it also follows from conditional negativity:
apply~\eqref{eq:CND} to frequencies $0,k$ and coefficients
$1,-1$ to obtain $-2\psi(k)\leq0$.

\section{The completed multiplier and its exact generator}
The analytic construction in this section requires only
$\psi(k)\geq0$. Conditional negativity is not needed.
Parseval's identity identifies $L^2(\T^3,m;\C)$ with
$\ell^2(\Lambda;\C)$. Since $0<e^{-t\psi(k)}\leq1$, its
diagonal multiplier defines a contraction $S_t$.
The identities $S_0=I$ and $S_sS_t=S_{s+t}$ hold coefficientwise.
For any fixed $u$ and $t_n\to t$ in $[0,\infty)$,
\[
 \norm{S_{t_n}u-S_tu}_2^2
 =\sum_k |e^{-t_n\psi(k)}-e^{-t\psi(k)}|^2|\widehat u(k)|^2
 \longrightarrow0.
\]
The summands are bounded by $4|\widehat u(k)|^2$; a finite-head,
small-tail argument proves the convergence, or equivalently
dominated convergence on counting measure. The multiplier is real,
so $S_t$ is self-adjoint. If $\psi(0)=0$, then $S_t1=1$.

\begin{proposition}\label{prop:generator}
For every nonnegative symbol, the strong right derivative
$Au=\lim_{t\downarrow0}(S_tu-u)/t$ exists exactly on
\eqref{eq:domain}. This operator is closed and has
$\mathcal P$ as a graph core.
\end{proposition}
\begin{proof}
If the strong derivative exists with value $v$, each continuous
Fourier coordinate functional gives
\[
 \widehat v(k)=\lim_{t\downarrow0}
       \frac{e^{-t\psi(k)}-1}{t}\widehat u(k)
             =-\psi(k)\widehat u(k).
\]
Parseval then forces $\sum_k\psi(k)^2|\widehat u(k)|^2<\infty$.
Conversely, if this sum is finite, let $v$ have those coefficients.
The scalar estimate $0\leq1-e^{-a}\leq a$ for $a\geq0$ gives
\[
 \left|\frac{e^{-t\psi(k)}-1}{t}+\psi(k)\right|
                                      \leq\psi(k),\qquad t>0.
\]
The scalar expression tends to zero for every $k$, so domination
by the summable squared sequence
$\psi(k)^2|\widehat u(k)|^2$ proves convergence to $v$ in $L^2$.

To prove closedness, suppose $u_n\to u$ and $Au_n\to v$ in
$L^2$. Coordinatewise limits give
$\widehat v(k)=-\psi(k)\widehat u(k)$, and since $v\in L^2$
the same characterization puts $u\in\Dom(A)$ with $Au=v$.
Finally truncate the Fourier series of $u\in\Dom(A)$ to
$\{k:|k_j|\leq N,\ j=1,2,3\}$. Both the $L^2$ tail and its
$A$-image tail tend to zero by their respective square sums.
These truncations are finite Fourier polynomials, proving the
core assertion in the norm $\norm u_2+\norm{Au}_2$.
\end{proof}

The semigroup can therefore exist with a completely controlled
closed generator even when it fails positivity. Section
\ref{sec:threshold} gives an explicit failure within this same
construction. It cannot be repaired by calling the coefficient
evolution a Markov operator before checking the physical cone.


\section{A bounded contraction that fails positivity at every positive time}\label{sec:bounded-rate}
A nonnegative, even, bounded rate can satisfy all elementary Hilbert
contraction requirements and still fail the positivity criterion.
Set $w=(1,1,1)\in\Z^3$ and define
\[
 \psi(k)=\begin{cases}1,&k=2w\text{ or }k=-2w,\\0,&\text{otherwise}.
 \end{cases}
\]
Thus $0\leq\psi\leq1$, $\psi(-k)=\psi(k)$ and $\psi(0)=0$.
Let $P$ be the Haar-$L^2$ orthogonal projection onto
$\operatorname{span}\{e_{2w},e_{-2w}\}$. The actual multiplier is
\begin{equation}\label{eq:boundedsemigroup}
 S_t=I+(e^{-t}-1)P,\qquad A=-P,\qquad\Dom(A)=L^2.
\end{equation}
On $\operatorname{ran}P$ its multiplier is $e^{-t}$, and on
$\ker P$ it is one. This proves the semigroup law, self-adjointness,
contraction and preservation of constants. It also proves operator-norm
continuity: $\|S_t-I\|=1-e^{-t}$ for $t\geq0$, and
$(S_t-I)/t\to-P$ in operator norm. Hence the entire space is the
closed generator domain; no unbounded-domain defect is responsible
for the failure below. The nonnegative Hilbert energy is
\[
 -\langle Au,u\rangle=\|Pu\|_2^2\geq0.
\]

The two-frequency covariance test is explicit:
\[
 \bigl(K_\psi(k,l)\bigr)_{k,l\in\{w,2w\}}
 =\begin{pmatrix}0&1/2\\1/2&1\end{pmatrix}.
\]
Indeed $\psi(w)=0$, $\psi(2w)=1$ and $\psi(2w-w)=0$.
For $f=e_w-\tfrac12e_{2w}$, equation~\eqref{eq:gammaenergy} gives
\begin{equation}\label{eq:boundedgamma}
 \Gamma_1(f,f)(0)=-\tfrac14.
\end{equation}
So the rate is neither conditionally negative nor an anchored
covariance symbol, despite boundedness, symmetry and nonnegativity.

There is a direct physical cone witness at \emph{every} $t>0$.
Write $\theta=2\pi w\cdot x$ and take
\[
 h(x)=(1-\cos\theta)^2
 =\tfrac32-2\cos\theta+\tfrac12\cos(2\theta)\geq0.
\]
Only its final harmonic is damped by \eqref{eq:boundedsemigroup};
therefore
\[
 (S_th)(0)=\tfrac12(e^{-t}-1)<0\qquad(t>0).
\]
This is a finite Fourier polynomial, so its image is continuous and
strict negativity at the origin persists on an open neighborhood.
Every such neighborhood has positive Haar measure: it contains a
small coordinate box in a torus chart, whose product Lebesgue measure
is positive. Thus the witness violates the actual almost-everywhere
Haar-$L^2$ cone, not just a choice of pointwise representative.

The bounded rate and the value \eqref{eq:boundedgamma} are mapped
in \code{BoundedRateGammaObstruction}. The whole-space projection
formula and the all-positive-time nonnegative test $h$ are explicit
human deductions from that rate and the already established physical
multiplier. No separate named formal theorem for $h$ is claimed.
This is a concrete instance of the classical positivity equivalence,
not a new general obstruction principle.

\section{Integrated Gamma energy can decay without pointwise positivity}
The distinction also survives exact integrated energy identities.
Fix $\gamma>0$, distinct integers $n_1,\ldots,n_N$ and complex
coefficients $c_j$. On the first coordinate circle use
\[
 \psi_\gamma(k)=|k_1|^\gamma,\qquad
 f(x)=\sum_{j=1}^N c_j e_{n_j e_1}(x),\qquad
 E_q(f)=\int_{\T}\Gamma_q(f,f)(x_1,0,0)\,dx_1.
\]
Here $q$ is a nonnegative integer. The zeroth iterate is $|f|^2$.
The integral is real, and orthogonality of distinct circle characters
in \eqref{eq:gammaenergy} gives the exact formula
\begin{equation}\label{eq:integratedgamma}
 E_q(f)=\sum_j |n_j|^{\gamma q}|c_j|^2\geq0,
\end{equation}
where the weight is interpreted as one when $q=0$, including
$n_j=0$. For $q\geq1$, the zero frequency has zero weight.
Indeed the integral eliminates every off-diagonal term and
$K_{\psi_\gamma}(n e_1,n e_1)=|n|^\gamma$.
This computation needs neither conditional negativity nor positivity
of the pointwise Gamma kernel.

Let $f_t=\sum_j e^{-t|n_j|^\gamma}c_j e_{n_j e_1}$. Then
\begin{align}\label{eq:energydecay}
 E_q(f_t)&=\sum_j |n_j|^{\gamma q}
                 e^{-2t|n_j|^\gamma}|c_j|^2,\\
 \frac{d}{dt}E_q(f_t)&=-2E_{q+1}(f_t).\nonumber
\end{align}
The derivative follows by differentiating a finite sum, so no closure
or infinite-series interchange has been used. If all frequencies
are nonzero, then $|n_j|^\gamma\geq1$ and
\begin{equation}\label{eq:integratedsharpdecay}
 E_q(f_t)\leq e^{-2t}E_q(f)\qquad(q\geq0,\ t\geq0).
\end{equation}
The single unit frequency attains equality for every $q$ and $t$,
so the coefficient two in the exponent is sharp at every energy
level. For $q\geq1$ a zero mode is harmless; for $q=0$ it must be
removed to obtain this decay for nonzero constants. Equivalently,
one can apply the zeroth-level estimate to the mean-zero part.

These identities hold for every positive $\gamma$, even
$\gamma>2$, when the pointwise positivity criterion fails by
Section~\ref{sec:threshold}. Thus nonnegative circle-integrated
energies, exact dissipation and sharp spectral decay do not prove
pointwise Gamma positivity, Markov positivity or a probability law.
For $0<\gamma\leq2$, the positive-law results below apply in
addition; for larger orders only the stated multiplier and finite
energy assertions are made. No completed pointwise Gamma domain is
inferred in either case.

The formal module \code{StableGammaEnergyDecay} records the actual
circle average, diagonal energy, finite semigroup, derivative and
sharp unit-frequency witness. Its nonzero-frequency decay statement
retains the nonzero-frequency hypothesis even at positive levels;
the harmless-zero-mode variant just proved is a labelled elementary
human extension of the same finite sum. These are classical Fourier
orthogonality and multiplier calculations. They supplement the
positivity equivalence rather than supply a new curvature theorem.

\section{From finite squares to the full Haar cone}
First calculate on a finite polynomial $f=\sum_i c_i e_{k_i}$:
\begin{equation}\label{eq:squareheat}
 (S_t|f|^2)(x)=\sum_{i,j}
      \overline{c_i e_{k_i}(x)}c_j e_{k_j}(x)
                         e^{-t\psi(k_j-k_i)}.
\end{equation}
This is equality of continuous finite Fourier representatives.
If the matrices~\eqref{eq:heatmatrix} are positive semidefinite,
the right side is real and nonnegative at every $x$.

\begin{lemma}\label{lem:cone}
The $L^2$ closure of $\{|f|^2:f\in\mathcal P\}$ is exactly
$\mathcal C$.
\end{lemma}
\begin{proof}
For $z\in\C$ set $\pi(z)=\max(\operatorname{Re}z,0)$, regarded
as a complex number with zero imaginary part. The map $\pi$ is
one-Lipschitz, vanishes at zero and induces a one-Lipschitz map
on $L^2$. Its fixed-point set is $\mathcal C$, hence this cone
is closed. Every polynomial square lies in it.

For the reverse inclusion, take $u\in\mathcal C$ and choose
continuous complex $v_n$ converging to $u$ in $L^2$. Such
approximation follows from continuous-function density for Haar
measure on the compact torus. The clipped continuous functions
$g_n=\pi(v_n)$ are nonnegative and converge to $u$, because
$\pi(u)=u$ almost everywhere and $\pi$ is one-Lipschitz.

For fixed $n$, the function $\sqrt{g_n}$ is continuous. Finite
Fourier polynomials are uniformly dense in continuous complex
functions on the torus, by the Stone--Weierstrass theorem applied
to the self-conjugate character algebra. Choose $p_j$ converging
uniformly to $\sqrt{g_n}$. Then
\[
 \norm{|p_j|^2-g_n}_\infty
 \leq\norm{p_j-\sqrt{g_n}}_\infty
              \bigl(\norm{p_j}_\infty+\norm{\sqrt{g_n}}_\infty\bigr)
 \longrightarrow0.
\]
Since $m$ is a probability measure, this implies $L^2$
convergence. Choose one such square within $1/n$ of $g_n$ in
$L^2$. The chosen squares converge to $u$.
\end{proof}

\begin{proposition}\label{prop:physical}
Under the hypotheses of Theorem~\ref{thm:main}, positivity of all
finite heat matrices is equivalent to $S_t\mathcal C\subseteq
\mathcal C$ for every $t\geq0$.
\end{proposition}
\begin{proof}
In the forward direction use~\eqref{eq:squareheat}, then approximate
$u\in\mathcal C$ by squares using Lemma~\ref{lem:cone}.
Boundedness of $S_t$ gives convergence of their images, and
closedness of $\mathcal C$ gives $S_tu\in\mathcal C$.

Conversely, apply physical positivity to every $|f|^2$.
Its image is a continuous finite polynomial, so almost-everywhere
nonnegativity implies pointwise nonnegativity. Indeed, a negative
value of its real part would give a nonempty open set of negative
values; Haar measure on the torus assigns positive measure to
every nonempty open set. Its imaginary part similarly vanishes
everywhere if it vanishes almost everywhere. Evaluating
\eqref{eq:squareheat} at $x=0$ and allowing arbitrary real $c_i$
proves nonnegativity of every heat-matrix quadratic form.
Evenness already supplies symmetry.
\end{proof}

\begin{proof}[Completion of the proof of Theorem~\ref{thm:main}]
Lemma~\ref{lem:anchor}, Proposition~\ref{prop:gammapositive},
Proposition~\ref{prop:schoenberg} and
Proposition~\ref{prop:physical} give all five equivalences and
the finite Gamma hierarchy. The semigroup properties and
Proposition~\ref{prop:generator} give the analytic assertions.
\end{proof}

For real bounded $u$ with $0\leq u\leq1$, positivity and
$S_t1=1$ also give $0\leq S_tu\leq1$ almost everywhere. Thus the
word Markov, when used for these operators, refers to this order
and constant-preservation property on the physical function space.
No sample-path regularity or filtration is built by the proof.

\section{The sharp active-coordinate threshold}\label{sec:threshold}
Let $a_j\geq0$ and $\alpha_j>0$ for $j=1,2,3$, and set
\begin{equation}\label{eq:mixed}
 \psi(k)=\sum_{j=1}^3 a_j|k_j|^{\alpha_j},\qquad 0^\alpha=0
       \quad(\alpha>0).
\end{equation}
A coordinate is active when $a_j>0$. The generator is defined
by these rates, with no alternative Euclidean radial power
substituted for their coordinatewise sum.

\begin{theorem}\label{thm:threshold}
For the symbol~\eqref{eq:mixed}, every condition in
Theorem~\ref{thm:main} holds if and only if
\begin{equation}\label{eq:active}
 a_j>0\ \Longrightarrow\ \alpha_j\leq2
               \quad(j=1,2,3).
\end{equation}
If an active exponent exceeds two, a nonnegative finite Fourier
square has a negative evolved value on a set of positive Haar
measure at some positive time.
\end{theorem}
\begin{proof}
We first prove the positive cases, including the endpoint.
For $0<\alpha<2$ put
\[
 I_\alpha=\int_0^\infty(1-\cos s)s^{-1-\alpha}\,ds.
\]
This is finite and strictly positive. Near zero use
$0\leq1-\cos s\leq s^2/2$; near infinity use $1-\cos s\leq2$;
positivity follows on an interval where the cosine is not one.
Changing variables gives, also at $x=0$,
\begin{equation}\label{eq:stableintegral}
 |x|^\alpha=I_\alpha^{-1}\int_0^\infty
                  (1-\cos(sx))s^{-1-\alpha}\,ds.
\end{equation}
For a real sample $x_i$ and coefficients summing to zero,
\[
 \sum_{i,l}c_i c_l(1-\cos(s(x_l-x_i)))
              =-\left|\sum_i c_i e^{isx_i}\right|^2\leq0.
\]
Every individual integrand in~\eqref{eq:stableintegral} is
integrable, so finite summation and integration commute. This
proves conditional negativity of $|x-y|^\alpha$.
At $\alpha=2$, expansion gives directly
\[
 \sum_{i,l}c_i c_l(x_l-x_i)^2
                   =-2\left(\sum_i c_i x_i\right)^2\leq0.
\]
Restrict to each integer coordinate, multiply by $a_j\geq0$
and add. If $a_j=0$ its summand is identically zero regardless
of the positive exponent. Thus~\eqref{eq:active} implies
conditional negativity of~\eqref{eq:mixed}, and the main
theorem supplies every asserted positivity.

For necessity, fix an active coordinate $j$ and sample
$0,e_j,2e_j$ with coefficients $(1,-2,1)$. Their sum is zero.
Directly summing the six off-diagonal terms gives
\begin{equation}\label{eq:triple}
 \sum_{i,l=0}^2 c_i c_l\psi((l-i)e_j)
                         =2a_j(2^{\alpha_j}-4).
\end{equation}
If $\alpha_j>2$, this is strictly positive, contradicting
conditional negativity.

The same witness gives a physical counterexample without any
appeal to an abstract converse. Let
\[
 p(x)=1-2e_{e_j}(x)+e_{2e_j}(x).
\]
Then $|p|^2=6-8\cos(2\pi x_j)+2\cos(4\pi x_j)\geq0$, and
\begin{equation}\label{eq:squarewitness}
 (S_t|p|^2)(0)=6-8e^{-a_jt}+2e^{-a_j2^{\alpha_j}t}.
\end{equation}
This vanishes at $t=0$ and has derivative
$-2a_j(2^{\alpha_j}-4)<0$ there. It is therefore negative for
all sufficiently small $t>0$. Each evolved polynomial is
continuous, so its negative set contains an open neighborhood
of zero and has positive Haar measure.
\end{proof}

This proof also shows exactly why a zero-weight exponent must
not be tested. For example $(a_1,a_2,a_3)=(1,0,0)$ and
$(\alpha_1,\alpha_2,\alpha_3)=(2,4,10)$ give the ordinary
one-coordinate quadratic heat evolution, with stationary modes
in the other two directions. Its positivity is unaffected by
the nominal inactive exponents. The all-zero symbol gives the
identity semigroup and is included as well.

The integral representation~\eqref{eq:stableintegral} is a
classical stable-symbol argument. In the formal development,
fractional covariance matrices provide the positive interior
case; the human proof supplies an equivalent direct integral
argument. No new formal declaration for that reorganization
is claimed.

\section{Finite-time increment kernels}
Assume from now on that $\psi$ satisfies Theorem~\ref{thm:main}.
For each $t\geq0$ define
\begin{equation}\label{eq:q}
 q_t(k)=1-e^{-t\psi(k)}.
\end{equation}
This symbol is normalized, even and nonnegative. It should be
distinguished from the original generator symbol $\psi$.

\begin{proposition}\label{prop:qt}
The anchored kernel $C_t=K_{q_t}$ is positive semidefinite at
every $t\geq0$, and
\begin{equation}\label{eq:incrementkernel}
 2C_t(k,l)=e^{-t\psi(l-k)}-e^{-t\psi(k)}-e^{-t\psi(l)}+1.
\end{equation}
If a torus-valued random variable $X_t$ has Fourier law
$\mathbb E e_k(X_t)=e^{-t\psi(k)}$, then
\begin{equation}\label{eq:randomincrement}
 2C_t(k,l)=\mathbb E\bigl[
       \overline{(e_k(X_t)-1)}(e_l(X_t)-1)\bigr].
\end{equation}
\end{proposition}
\begin{proof}
For any zero-sum real $c_i$, the constant part of $q_t$ cancels,
so
\[
 \sum_{i,j}c_i c_jq_t(k_j-k_i)
              =-\sum_{i,j}c_i c_je^{-t\psi(k_j-k_i)}\leq0
\]
by the finite heat-matrix positivity. Lemma~\ref{lem:anchor}
gives positivity of $C_t$. Expanding its definition proves
\eqref{eq:incrementkernel}. Expanding the product of character
increments gives~\eqref{eq:randomincrement} under the stated
Fourier law, using $e_{-k}e_l=e_{l-k}$ and evenness.
\end{proof}

The expectation in~\eqref{eq:randomincrement} is an uncentered
second-moment Gram kernel of the increments $e_k(X_t)-1$.
Those increments need not have mean zero. The word covariance
in ``anchored covariance kernel'' denotes the positive-kernel
structure; it must not silently turn this identity into the
centered covariance of those random variables. The frozen library
constructs actual transition laws and fixed-step chains with
this Fourier law. The proposition's matrix proof does not depend
on a continuous-time path construction, and none is used here.

As $t\downarrow0$, $q_t(k)/t\to\psi(k)$ for each frequency.
Hence $C_t(k,l)/t\to K_\psi(k,l)$ on each finite sample.
This is finite-dimensional entrywise convergence, not an
operator-norm assertion over all frequencies.

\section{Positivity does not imply increase in time}
The scalar values $q_t(k)$ increase in $t$ for every fixed
$k$. Their anchored matrices need not increase in Loewner
order. Recall that $C_t\succeq C_s$ means the difference is
positive semidefinite on the chosen common finite sample.

\begin{proposition}\label{prop:nonmonotone}
Take $\psi(k)=k_1^2$, frequencies $k=e_1$ and $l=-2e_1$,
and the real vector $b=(1,-2)$. For $s=\log2$ and $t=2\log2$,
\begin{equation}\label{eq:negative}
 b^{\mathsf T}(C_t-C_s)b=-\frac{16897}{131072}<0.
\end{equation}
Both $C_s$ and $C_t$ are positive semidefinite.
\end{proposition}
\begin{proof}
The rates of $k,l,l-k$ are $1,4,9$. Writing $r=e^{-u}$ gives
\[
 C_u(k,k)=1-r,\qquad C_u(l,l)=1-r^4,
 \qquad C_u(k,l)=\frac{1-r-r^4+r^9}{2}.
\]
Consequently
\[
 b^{\mathsf T}C_ub=3+r-2r^4-2r^9.
\]
At $u=s$ one has $r=1/2$; at $u=t$ one has $r=1/4$.
Subtracting gives
\[
 \left(\frac14-\frac12\right)
 -2\left(\frac1{256}-\frac1{16}\right)
 -2\left(\frac1{262144}-\frac1{512}\right)
                         =-\frac{16897}{131072}.
\]
The quadratic symbol meets~\eqref{eq:active}, so each individual
matrix is positive by Proposition~\ref{prop:qt}.
\end{proof}

The example distinguishes two assertions: positivity at each
time and positivity of time differences. It does not contradict
a curvature theorem whose hypotheses and conclusion are
different from Loewner monotonicity of $C_t$.

There is a similar distinction within the iterated Gamma
hierarchy. For $\psi(k)=\tfrac12|k_1|$ and $f=e_{e_1}$,
\eqref{eq:modegamma} gives
\[
 \Gamma_n(f,f)(x)=2^{-n}.
\]
Thus $\Gamma_2(f,f)=1/4<1/2=\Gamma_1(f,f)$ despite positivity
of every iterate. No inequality $\Gamma_2\geq\Gamma_1$ follows
from the hierarchy alone. The classical Bakry--\'Emery
calculus~\cite{bakryemery} asks for specified quantitative
comparisons, with their normalization and domain hypotheses;
the jump-type power symbols here also need not satisfy the
diffusion chain rule. The finite algebraic result must retain
those distinctions.

\section{Quadratic flatness on the finite core}
For the same one-coordinate quadratic symbol, the precise
pointwise curvature statement has sharp constant zero. This
conclusion has its own proof; it is not inferred from
Proposition~\ref{prop:nonmonotone}.

\begin{proposition}\label{prop:flat}
For $\psi(k)=k_1^2$, one has $\Gamma_2(f,f)(x)\geq0$ for every
$f\in\mathcal P$ and every $x$. There is no $\rho>0$ such that
\begin{equation}\label{eq:curvature}
 \Gamma_2(f,f)(x)\geq\rho\Gamma_1(f,f)(x)
\end{equation}
for all real-valued $f\in\mathcal P$ and all $x\in\T^3$.
\end{proposition}
\begin{proof}
The first assertion follows from the main theorem, or directly
from $K_\psi(k,l)=k_1l_1$ in~\eqref{eq:gammaenergy}. To test
\eqref{eq:curvature}, let
\[
 f=ie_{e_1}-ie_{-e_1}=-2\sin(2\pi x_1).
\]
This polynomial is real-valued. In its two-mode presentation
the frequencies are $e_1,-e_1$, the coefficients are $i,-i$,
and at the origin the first kernel matrix is
$\left(\begin{smallmatrix}1&-1\\-1&1\end{smallmatrix}\right)$.
Its energy is $|i-(-i)|^2=4$. The entrywise square matrix is
the all-ones matrix, whose energy is $|i+(-i)|^2=0$.
Thus $\Gamma_1(f,f)(0)=4$ and $\Gamma_2(f,f)(0)=0$.
No positive $\rho$ can satisfy~\eqref{eq:curvature}.
\end{proof}

For a check on normalization, set $D=(2\pi)^{-1}\partial_{x_1}$
on $\mathcal P$. Then $L_\psi=D^2$,
$\Gamma_1(f,f)=|Df|^2$ and $\Gamma_2(f,f)=|D^2f|^2$.
These identities follow either by differentiation or from the
same mode formula. They reproduce the sine values above.
Positivity on the physical cone therefore coexists with sharp
zero pointwise curvature in this example. The finite core is
the domain of this statement; no extension to arbitrary
members of $\Dom(A)$ is hidden in its notation.


\section{Actual probability laws and fixed-step chains}\label{sec:physical-laws}
The positive multipliers have an actual measure realization, not just
an abstract coefficient interpretation. We retain the hypotheses of the
main equivalence: $r:\Lambda\to[0,\infty)$ is even, $r(0)=0$, and its
anchored kernel is positive semidefinite. Write $P_t e_k=e^{-tr(k)}e_k$.

\begin{theorem}\label{thm:laws}
For each $t\geq0$ there is a unique probability measure $\mu_t$ on
$\T^3$ with
\[
 \int e_{-k}(z)\,d\mu_t(z)=e^{-tr(k)}\qquad(k\in\Lambda).
\]
These measures satisfy $\mu_0=\delta_0$ and
$\mu_s*\mu_t=\mu_{s+t}$, and depend weakly continuously on $t$.
The translated laws $Q_t(x,A)=\mu_t(A-x)$ are Markov kernels satisfying
Chapman--Kolmogorov. On continuous functions and on Haar-$L^2$ their
integral operators agree with the original Fourier multipliers.
\end{theorem}
\begin{proof}
Let $B_N=\{-N,\ldots,N\}^3$ and
\[
 q_N=|B_N|^{-1/2}\sum_{j\in B_N}e_j,\qquad h_N=|q_N|^2.
\]
Orthogonality gives $\int h_N\,dm=1$. By finite-square positivity,
$g_{N,t}=P_t h_N$ is a continuous nonnegative polynomial of integral
one. Thus $g_{N,t}m$ is a probability measure. Its $k$th Fourier moment is
\[
 e^{-tr(k)}a_N(k),\qquad
 a_N(k)=\frac{|\{j\in B_N:j+k\in B_N\}|}{|B_N|}.
\]
For $M=\max_i|k_i|$ and $N\geq M$, the numerator lies between
$|B_{N-M}|$ and $|B_N|$. Therefore $a_N(k)\to1$ for every fixed $k$.
Probability measures on the compact metric torus have a weakly convergent
subsequence. Integrating each bounded continuous character along that
subsequence gives all the required moments of its limit.

If two probability measures have these moments, their integrals agree
on trigonometric polynomials. These form a uniformly dense subalgebra
of $C(\T^3;\C)$: they contain constants, are closed under conjugation
and separate points. Uniform approximation and the bound
$|\int f\,d\mu|\leq\|f\|_\infty$ extend equality to every continuous
function. Continuous functions determine finite measures, proving
uniqueness. At time zero $\delta_0$ has the prescribed moments.
Fubini and $e_k(x+y)=e_k(x)e_k(y)$ show that the moments of
$\mu_s*\mu_t$ multiply to those of $\mu_{s+t}$, so uniqueness proves
the convolution identity.

If $t_j\to t$, every weakly convergent subsequence of $\mu_{t_j}$
has the moments of $\mu_t$, by continuity of $e^{-tr(k)}$.
Compactness and uniqueness force the full sequence to converge to
$\mu_t$. This proves weak continuity. For any Borel set $A$,
$Q_t(x,A)=\int1_A(x+z)\,d\mu_t(z)$ is Borel in $x$ by measurability
of addition and integration of a bounded Borel function. It has total
mass one. Fubini and convolution give
\[
 \int Q_t(y,A)Q_s(x,dy)=Q_{s+t}(x,A).
\]
Its action on $e_k$ is $e^{-tr(k)}e_k(x)$, since evenness gives the
same real moment at $k$ and $-k$. The action on continuous functions is
a contraction in the uniform norm. Moreover Jensen and Haar invariance
give $\int|Q_tf|^2dm\leq\int|f|^2dm$. Density of polynomials in each
space identifies these integral operators with the respective
multiplier extensions. This proves the last assertion on the original
physical Haar space, not just on a coefficient space.
\end{proof}

Fix a step $h\geq0$. On the countable product probability space with
independent coordinates $Z_j$ of law $\mu_h$, set
$X_0=x$ and $X_n=x+Z_1+\cdots+Z_n$. This is a human reorganization of
the source's Ionescu--Tulcea construction. Independence gives, first
for bounded Borel $f$ and then as a regular conditional law,
\[
 \mathbb E[f(X_{n+1})\mid X_0,\ldots,X_n]=Q_hf(X_n).
\]
Indeed the finite past is measurable in $Z_1,\ldots,Z_n$ and $Z_{n+1}$
is independent of that larger sigma field; conditioning the resulting
identity again gives the displayed formula. Convolution shows
$\mathcal L(X_n)=Q_{nh}(x,\cdot)$, including $h=0$ and $n=0$.
The finite-dimensional probabilities are the iterated kernels
$\delta_x(dx_0)\prod_{j=1}^nQ_h(x_{j-1},dx_j)$, which also identify
the resulting countable-path law with the source trajectory law.
No sample paths at non-grid times have been constructed here.

For $x=0$, the actual first-step character increment pairing is
\begin{equation}\label{eq:actual-pair}
 \mathbb E[(e_{-k}(X_1)-1)(e_l(X_1)-1)]
 =e^{-hr(k-l)}-e^{-hr(k)}-e^{-hr(l)}+1.
\end{equation}
Expand the product and use its three character moments to prove this
identity. Every integrand is bounded, so no interchange requires an
extra domain. Its right derivative at $h=0$ is
$r(k)+r(l)-r(k-l)=2K_r(k,l)$. Varying $h$ varies the law of the
first-step chain; this derivative is not a claim of differentiable
sample paths or of an unsmoothed pointwise Gamma domain.

\section{Poisson mixtures, an atom and a jump covariance}\label{sec:poisson-law}
Fix an original heat step $\delta\geq0$ and a finite Poisson parameter
$\lambda\geq0$. Define an actual probability measure by
\begin{equation}\label{eq:compound}
 \nu_\lambda=e^{-\lambda}\sum_{n=0}^\infty
        \frac{\lambda^n}{n!}\mu_{n\delta}.
\end{equation}
At $\lambda=0$ use $0^0=1$ in the $n=0$ term. The weights are
nonnegative and sum to one, so the sum defines a probability measure.
This elementary Poisson subordination is classical; see Varadhan's
primary notes, Sections 1.2 and 2.2~\cite{varadhan}.

\begin{theorem}\label{thm:compound}
Put $q_\delta(k)=1-e^{-\delta r(k)}$. Then
\[
 \widehat\nu_\lambda(k)=e^{-\lambda q_\delta(k)},\qquad
 \nu_\lambda*\nu_\eta=\nu_{\lambda+\eta},\qquad
 \nu_0=\delta_0.
\]
The transformed symbol $q_\delta$ is normalized, even, nonnegative
and has positive semidefinite anchored kernel. Its unique physical
heat law at parameter $\lambda$ is exactly $\nu_\lambda$.
Nevertheless
\begin{equation}\label{eq:atom}
 \nu_\lambda\geq e^{-\lambda}\delta_0,
 \qquad \nu_\lambda\not\ll m
 \quad\text{for every finite }\lambda\geq0.
\end{equation}
\end{theorem}
\begin{proof}
Characters have modulus one, so the sum of their absolute integrals
weighted by the Poisson probabilities is at most one. Thus integration
commutes with the sum in~\eqref{eq:compound}, and
\[
 \widehat\nu_\lambda(k)
 =e^{-\lambda}\sum_{n\geq0}\frac{(\lambda e^{-\delta r(k)})^n}{n!}
 =\exp[-\lambda(1-e^{-\delta r(k)})].
\]
Multiplication of these moments and Fourier uniqueness prove the
convolution assertion and the zero-parameter assertion.
Normalization, evenness and nonnegativity of $q_\delta$ follow from
those of $r$. Formula~\eqref{eq:actual-pair} for one original heat
jump $Z$ of law $\mu_\delta$ gives
\begin{equation}\label{eq:jump-gram}
 \mathbb E[(e_{-k}(Z)-1)(e_l(Z)-1)]
 =q_\delta(k)+q_\delta(l)-q_\delta(k-l)
 =2K_{q_\delta}(k,l).
\end{equation}
For finitely many coefficients $c_j$, multiplying by
$\overline{c_i}c_j$ and summing gives the nonnegative integral
$\mathbb E|\sum_jc_j(e_{k_j}(Z)-1)|^2$. Thus the anchored kernel
is positive semidefinite. Theorem~\ref{thm:laws} applies to
$q_\delta$; its moments and uniqueness identify its law with
\eqref{eq:compound}.

The $n=0$ term in that sum is $e^{-\lambda}\mu_0
=e^{-\lambda}\delta_0$; every remaining term is a nonnegative
measure. This proves the measure inequality, not merely positivity
of one coefficient. Haar measure has null singletons: if a singleton
had mass $a>0$, translation invariance would give the same mass to
each of $N$ distinct points, forcing $Na\leq1$ for every $N$.
Hence $m\{0\}=0$ while $\nu_\lambda\{0\}\geq e^{-\lambda}>0$.
This proves failure of absolute continuity. The mass at zero may be
larger than $e^{-\lambda}$; no equality is asserted. In particular
$\delta=0$ yields $\nu_\lambda=\delta_0$ at all parameters.
\end{proof}

The source atom theorem assumes that the reference measure has a null
origin singleton. The preceding Haar argument explicitly discharges
that hypothesis here; it is a classical deduction, not an additional
named Lean declaration. Thus actual positivity, conservativity and
weak continuity coexist with failure of absolute-continuity smoothing.

There is also a bounded operator explanation. Let $U=P_\delta$ on
$C(\T^3)$ or Haar-$L^2$. Since $\|U\|\leq1$, the series
\[
 V_\lambda=e^{-\lambda}\sum_{n\geq0}\frac{\lambda^n}{n!}U^n
             =\exp[\lambda(U-I)]
\]
converges in operator norm. The norm tail is bounded by the scalar
exponential tail. The convolution identity gives $U^n=P_{n\delta}$,
so $V_\lambda$ is exactly integration against translated
$\nu_\lambda$, initially on bounded continuous functions and then on
$L^2$ by contraction and density. If $A=U-I$, then $\|A\|\leq2$ and
\[
 \left\|\frac{e^{\lambda A}-I}{\lambda}-A\right\|
 \leq\sum_{j\geq2}\frac{\lambda^{j-1}\|A\|^j}{j!}\longrightarrow0.
\]
Consequently the generator in the parameter $\lambda$ is the bounded
operator $P_\delta-I$ on the whole space. This is an explicitly proved
classical consequence, separate from the original unbounded generator
domain for $r$. Bounded jump generation does not imply smoothing.

For each fixed $\lambda$, independent increments of law $\nu_\lambda$
give a countable-step chain exactly as above. Its first step from zero
has the atom~\eqref{eq:atom}, and its character increment pairing is
\[
 e^{-\lambda q_\delta(k-l)}-e^{-\lambda q_\delta(k)}
       -e^{-\lambda q_\delta(l)}+1.
\]
The right derivative at parameter zero is
$2K_{q_\delta}(k,l)$, which by~\eqref{eq:jump-gram} is the second
moment of the increments of \emph{one original heat jump}. It is
an uncentered Gram pairing. A centered covariance would additionally
subtract the product of the increment means. This changes the
finite-parameter value, although that product has derivative zero at
zero. The derivative varies the one-step Poisson parameter across a
family of discrete chains; it does not construct a continuous-time
compound-Poisson process with cadlag paths.

\section{Formal source correspondence and editorial scope}
The mathematical source is frozen at
\begin{center}\small\ttfamily
bab688bb6140cd6316c80680095a914ea0d346db.
\end{center}
The namespace prefix is
\begin{center}\code{OperatorDerivativeRewrite.Foundation}\end{center}
The accompanying \code{theorem-map.json} records full
declaration names, exact lines, source hashes and baseline links.
The main correspondences are as follows.

\begin{longtable}{@{}p{0.27\textwidth}p{0.67\textwidth}@{}}
\textit{Paper statement}&\textit{Lean namespace and declaration}\\\hline
Anchoring&\code{FourierCovarianceCriterion.symbol_iff}\\[3pt]
Finite Gamma hierarchy&\code{FourierGammaCovariance.gamma_nonnegative}\\[3pt]
Normalized first test&\code{FourierGammaNormalization.normalized_gamma_iff_negative_even}\\[3pt]
Complex reality&\code{FourierGammaComplexCriterion.gamma_full_iff_covariance}\\[3pt]
Finite heat criterion&\code{FiniteSchoenberg.heat_positive_iff_covariance}\\[3pt]
Exact physical generator&\code{PhysicalTorusWaveSemigroup.generator_domain}\\[3pt]
Closed graph&\code{PhysicalTorusWaveSemigroup.graph_closed}\\[3pt]
Full cone closure&\code{PhysicalHeatPositiveCone.positive_l2_iff_square_closure}\\[3pt]
Full cone preservation&\code{PhysicalHeatPositiveCone.physical_evolution_positive_all}\\[3pt]
Mixed-symbol criterion&\code{AnisotropicMixedThreshold.anisotropic_negative_iff_active}\\[3pt]
Physical sharp threshold&\code{PhysicalHeatPositiveCone.anisotropic_physical_positive_iff_active}\\[3pt]
Finite-time kernel&\code{PhysicalHeatChain.heatIncrementSymbol_covariance}\\[3pt]
Actual increment pairing&\code{PhysicalHeatChain.chainLaw_increment_pair_gamma}\\[3pt]
Exact negative difference&\code{PhysicalHeatCovarianceNonmonotonicity.covariance_increment_energy_eq}\\
Actual physical laws&\code{PhysicalHeatFiniteBox.exists_physical_heat_law}\\[3pt]
Poisson mixture and atom&\code{PhysicalHeatCompoundPoisson.compoundHeatLaw_character} and \code{PhysicalHeatCompoundPoissonAtom.noJump_dirac_le}\\[3pt]
Sharp zero core curvature&\code{PhysicalHeatFlatCurvature.no_positive_quad_curvature}\\
\end{longtable}

The finite-polynomial graph-core assertion is a classical
consequence of the exact domain, proved by truncation in
Proposition~\ref{prop:generator}; no separate three-dimensional
physical graph-core declaration is claimed in the source map.
The explicit square calculation also shows failure throughout
a sufficiently small right neighborhood of zero, whereas the
mapped physical counterexample declaration asks only for one
positive time. These elementary deductions are proved here,
not attributed to additional named Lean theorems.
All physical statements use the unit-period Haar torus. The baseline
receipt records a full build and 54,943 declarations audited with the
standard three axioms; its evidence hashes were checked. No new Lean
build or mathematical source change is claimed.

This paper is a coherent exposition of classical positivity
structure and its verified physical realization, with explicit
scope and counterexamples. Schoenberg's primary paper
\cite{schoenberg} is an antecedent for the exponential
positive-definiteness equivalence, not merely a related
formalization. The power threshold and integral argument are
classical. We do not claim novelty for those statements or for
Schur products. Priority for the exact finite witnesses and
this combined formalization package remains unresolved. The Poisson mixture, bounded jump generator and no-jump mechanism are classical; Varadhan's primary notes were checked at the characteristic-function and operator-series formulas. No priority is claimed for them. The countable independent-increment realization reorganizes the source's kernel trajectory proof without asserting continuous-time path regularity.

The complete $L^2$ generator domain does not identify a closed
unsmoothed pointwise Gamma domain. Gamma remains on the finite Fourier
algebra. No continuous-time path regularity, positive curvature
constant or physical gauge-field mass gap is established.

\clearpage
\begin{thebibliography}{9}
\bibitem{varadhan} S. R. S. Varadhan,
\emph{Poisson Processes}, Spring 2006 lecture notes, Sections 1.2 and 2.2.
\url{https://math.nyu.edu/faculty/varadhan/spring06/spring06.1.pdf}.

\bibitem{schoenberg} I. J. Schoenberg,
Metric spaces and positive definite functions,
\emph{Trans. Amer. Math. Soc.} \textbf{44} (1938), 522--536.
\href{https://doi.org/10.1090/S0002-9947-1938-1501980-0}
{doi:10.1090/S0002-9947-1938-1501980-0}.
\bibitem{bakryemery} D. Bakry and M. \'Emery,
Diffusions hypercontractives,
\emph{S\'eminaire de Probabilit\'es XIX},
Lecture Notes in Math. \textbf{1123}, Springer, 1985, 177--206.
\url{https://www.numdam.org/item/SPS_1985__19__177_0/}.
\end{thebibliography}
\end{document}
